\documentclass[a4paper,11pt]{ltjsarticle}

% --- フォント設定 ---
\usepackage{luatexja}
\usepackage[haranoaji]{luatexja-preset}

% --- レイアウト設定 ---
\usepackage[top=20mm, bottom=20mm, left=20mm, right=20mm]{geometry}

% --- 箇条書き調整 ---
\usepackage{enumitem}

% --- セクション見出し調整 ---
\usepackage{titlesec}
\titleformat{\section}{\large\bfseries}{\thesection}{1em}{}
\titlespacing*{\section}{0pt}{1.0ex plus .2ex minus .2ex}{0.5ex plus .2ex}

% --- 行間設定 ---
\renewcommand{\baselinestretch}{0.98}

\begin{document}

% --- ヘッダー ---
\begin{flushright}
2026年1月16日
\end{flushright}
\vspace{-1.5em}

% --- タイトル ---
\begin{center}
    {\Large \textbf{「はずかしそうな笑い」の深層}} \\[0.3em]
    {\large ―『きつねのお客さま』：隠し続けた「本能」と癒やされた孤独―}
\end{center}

\vspace{-0.5em}

% --- 氏名 ---
\begin{flushright}
 4321　野秋琳太郎
\end{flushright}

\vspace{1.0em} 

\section{はじめに}
　あまんきみこの『きつねのお客さま』は、捕食者であるキツネが、本来は獲物であるヒヨコたちを守るためにオオカミと戦い、命を落とす物語である。

この作品の結末について、多くの読者は「愛による改心」や「自己犠牲」という美しい側面に着目する\cite{miyagawa, yoshimura}。しかし、テクストを精密に読むと、キツネの最期には不可解な表現が残されている。
彼は「安らかに」でも「誇らしげに」でもなく、「はずかしそうに」笑って死んでいるのである。

なぜ、命を懸けて他者を守った英雄的な瞬間に、彼は「恥じらい」を感じたのか。
本稿では、この「はずかしそうな笑い」という表現に焦点を絞って考察する。
キツネが抱えていた、「食欲」という名の「隠し事」と、それを凌駕していった「孤独」の心理を読み解くことで、彼の死が単なる悲劇ではなく、過去の自分自身を乗り越えた「自己超克」の証であったことを明らかにしたい。

\section{「隠し事」の正体と、孤独による変容}
　考察の前提として、キツネの動機が二重構造になっている点を確認したい。当初、キツネは物理的な「空腹」を満たすためにヒヨコたちを家に招いた。この時点での彼の「隠し事」は、明確な悪意である。

\begin{quote}
きつねは心の中で、にやりとわらった。\\
「よしよし、おれのうちにきなよ。」（注1）
\end{quote}

この「にやりとわらった」という描写は、彼がヒヨコを純粋な「餌」として見ていた動かぬ証拠である。
しかし、ヒヨコたちが発した「やさしいおにいちゃん」という言葉は、キツネ自身も自覚していなかった精神的な「孤独」を刺激した。

キツネは、「食べるための隠し事（罠）」を抱えたまま、奇妙な共同生活を続けることになる。
ここで重要なのは、彼がヒヨコたちを受け入れた理由が、いつしか「食欲」から「孤独の解消」へとすり替わっていったことだ。
物理的な空腹を満たすことよりも、精神的な承認欲求が満たされる快感の方が上回ってしまったのである。この「動機の逆転」こそが、彼を悲劇的な英雄へと変える鍵であった。

\section{オオカミという「鏡」との対決}
　物語のクライマックスで現れるオオカミは、単なる外敵ではない。考察の視点を広げれば、オオカミは「かつてのキツネ自身」の投影として読むことができる。

オオカミは、ヒヨコたちを純粋な「餌」としてしか見ていない。それは、物語の冒頭で「にやり」と笑っていたキツネの姿そのものである。キツネがオオカミと戦う構図は、彼が「保護者という新しい自分」を守るために、「捕食者という過去の自分」と戦っている姿の隠喩（メタファー）となっている。

\begin{quote}
きつねのからだに、勇気が、りんりんとわいた。\\
おお、たたかったとも、たたかったとも。（注2）
\end{quote}

ここで湧いた「勇気」は、単なる生存本能ではない。彼がボロボロになりながらも一歩も引かなかったのは、ヒヨコを守るためであると同時に、自分自身が再び「ただの獣」に戻ってしまうことを拒絶するためだったのではないだろうか。

\section{「はずかしそうな笑い」に込められた二つの意味}
　オオカミを退け、息絶える瞬間に浮かべた「はずかしそうな笑い」。この複雑な表情は、直後に動物たちが捧げた言葉と対比させることで、その意味が浮かび上がってくる。

\begin{quote}
そして、世界一やさしい、親切な、かみさまみたいな、そのうえゆうかんなきつねのために、なみだをながしたとさ。（注3）
\end{quote}

動物たちはキツネを「神様みたい」と称え、涙を流す。
しかし、キツネの胸中には、最後まで「最初は食べるつもりだった」という真っ黒な「隠し事」が存在していた。
このギャップこそが、「恥じらい」の正体である。

この「恥じらい」には二つの意味が込められている。
第一に、その「隠し事」を墓場まで持っていけたことへの安堵だ。
もし生き残ってしまえば、いつかその邪心がバレるかもしれない。
しかし、今ここで死ぬことで、彼はその「捕食者としての本能」を永遠に封印し、ヒヨコたちの中で「完璧な優しいお兄ちゃん」のままでいられる。

第二に、分不相応な幸福に対する「照れ」である。
「俺はそんな神様みたいな立派なやつじゃないよ」という謙遜と、それでも孤独だった自分が最期に愛されたことへのどうしようもない喜び。
自分のような者が英雄として扱われる状況へのこそばゆさが、最期の「はずかしそうな笑い」(注4)に集約されている。

\section{結論}
　キツネの死は、肉体的には敗北であり、悲劇的な結末に見える。
しかし、その表情に注目することで、精神的な側面における彼の勝利が見えてくる。

彼は最期に、「食べる／食べられる」という残酷な自然のルールから抜け出し、他者と心を通わせる世界に到達した。
彼が守り抜いたのは、ヒヨコたちの命であると同時に、彼らとの間に築き上げた「信頼」という虚構であった。
その虚構は、命を懸けたことで真実へと昇華された。キツネは「隠し事」を抱えた詐欺師として、少しだけ恥じらいながら、しかし孤独から解放されて旅立ったのである。

% --- 注・参考文献 ---
\vspace{2em}

\section*{注}
\begin{enumerate}[nosep]
    \item[(注1)] 授業資料集 p.27 上段9行目
    \item[(注2)] 授業資料集 p.28 上段11行目
    \item[(注3)] 授業資料集 p.28 下段2行目
    \item[(注4)] 授業資料集 p.28 上段15行目
\end{enumerate}

% --- 参考文献環境 ---
\begin{thebibliography}{99}
    \bibitem{miyagawa} 宮川久美「小学校国語教材『きつねのおきゃくさま』の主題についての考察」『奈良佐保短期大学研究紀要』第24号、2016年、pp.1-8。
    \bibitem{yoshimura} 吉村真理子、上田和美「言葉による伝え合いから生まれる共感的理解 ―『ずーっとずっとだいすきだよ』『きつねのおきゃくさま』読み語りの実践報告―」『千葉敬愛短期大学紀要』第44号、2022年、pp.95-105。
\end{thebibliography}

\end{document}